\section*{Capítulo 11}

\solucion{11.1} Factor para $m=+0.10$: $0.10/1.10=0.09091$. Sin cobertura: $4{,}704.0/3.4379\times0.09091=1{,}368.3\times0.09091=124.4$. Solo CCS: $938.4/3.4379\times0.09091=272.96\times0.09091=24.8$. CCS + NDF: $187.7/3.4379\times0.09091=54.60\times0.09091=5.0$.

\solucion{11.2} Nocional del NDF: S/ 938.4 millones. PCG $=-4{,}704.0+3{,}765.6+938.4=0$. La pérdida cambiaria ante una caída de 10\,\% sería nula, salvo el costo del NDF (aproximadamente $938.4\times(1/3.4367-1/3.4347)=-0.16$ millones por trimestre). La política permite hasta 100\,\%; el informe elige 80\,\% porque el saldo de cuentas por pagar fluctúa y cubrir el 100\,\% podría dejar a la empresa sobrecubierta si el saldo baja.

\solucion{11.3} NDF: $750.7/3.30-750.7/3.4347=227.48-218.56=+8.92$ millones. Cuentas por pagar: $750.7/3.4367-750.7/3.30=218.44-227.48=-9.05$ millones. Neto: $-0.13$ millones, el costo de la cobertura.

\solucion{11.4} El préstamo del BN tiene un cronograma largo y (casi) conocido: un CCS amortizable lo replica pago por pago y elimina el riesgo durante toda su vida. Las cuentas por pagar rotan, cambian de monto y vencen en semanas: un instrumento de corto plazo y renovable, como el NDF, permite ajustar el nocional cada trimestre. Usar NDF para el préstamo obligaría a renovar durante más de dos años a forwards desconocidos; usar un CCS para las cuentas por pagar fijaría un cronograma que no existe.

\solucion{11.5} Que esperar es especular: nadie sabe hacia dónde irá el tipo de cambio, y el forward no lo predice. Mientras se espera, la empresa sigue expuesta a una pérdida de unos US\$ 152 millones si el TC cae 10\,\%. La política de cobertura debe prohibir decisiones basadas en pronósticos y fijar límites de cobertura independientes de la opinión sobre el mercado.

\solucion{11.6} Porque las dos tasas están en monedas distintas y no se comparan directamente. El 5.87\,\% (5.98\,\% con CVA) es lo que cuesta, en dólares, una deuda sintética con el perfil del préstamo del BN; la alternativa para endeudarse en dólares sería el mercado de bonos, que exige a PETROPERÚ alrededor de 9.7\,\%. Además, parte del 5.87\,\% refleja que el cupón del BN (5.55\,\%) está por encima de la curva soberana en soles: el préstamo ya era \enquote{caro} en soles.

\solucion{11.7} De enero a junio de 2027 solo se pagan intereses: $6\times16.99=101.9$. De julio a diciembre se pagan seis cuotas francesas: $6\times218.28=1{,}309.7$. Total: 1,411.6. En 2028 hay doce cuotas: $12\times218.28=2{,}619.4$, casi el doble, porque toda la amortización se concentra en los últimos 18 meses tras el período de gracia.
